\subsection{A smaller abelian gate for two-prime mixed exponents}
\label{subsec:lm-mixed-compact-gate}

Let $q\ne3$ be an odd prime, and let
\[
 n=3^r q^s,\qquad r,s\ge1,\qquad S=\{3,q\}.
\]
The boundary-localization range will be determined exactly below. It
includes every $3q^s$ and, for arbitrary $r$, all primes
$q\equiv2,4,5,7\pmod9$.
As before, write $X_n$ for the compact Fermat curve,
$D_n$ for its coordinate cusps, $U_n=X_n\setminus D_n$,
$J_n=\operatorname{Jac}(X_n)$, and
$N_n=\pi_1(U_{n,\overline{\mathbb Q}})^{\mathrm{pro}\text{-}3}$.
The statements below retain an explicit compact Selmer hypothesis.
It and the higher graded conditions are established for $n=5\cdot3^r$
in Theorem~\ref{thm:lm-five-power-selmer} and for $n=7\cdot3^r$
by Corollary~\ref{cor:lm-septic-class-gate} and
Proposition~\ref{prop:septic-three-class}. For other exponents
they remain explicit hypotheses.

\begin{lemma}[Exact order of a cusp difference]
\label{lem:lm-exact-cusp-order}
For every $n\ge3$, the difference of two distinct coordinate
cusps on $X_n$ has exact order $n$ in $J_n(\overline{\mathbb Q})$.
\end{lemma}
This is the classical cusp-order assertion recorded by
\citet[p.~191, immediately before Theorem~A]{Coleman1986}.

\begin{lemma}[The abelian unit localization input]
\label{lem:lm-abelian-leopoldt}
If $F/\mathbb Q$ is a finite abelian extension, then
\[
 \mathcal O_F^\times\widehat{\ }_3
 \longrightarrow
 \prod_{v\mid3}(F_v^\times)^{\wedge_3}
\]
is injective.
\end{lemma}
\begin{proof}
This is the abelian case of Leopoldt's theorem.
\citet[pp.~192--193, Remark~1 after Theorem~$1_p$]{Brumer1968}
records the nonvanishing regulator for real abelian fields;
the unit-index assertion on p.~195 gives the CM extension.
It follows from the Baker--Brumer theorem.  We recall the application
to make its precise strength clear.  Modulo roots of unity,
the rational representation of
$G=\operatorname{Gal}(F/\mathbb Q)$ on global units contains
each relevant nontrivial even character with multiplicity
one, by Dirichlet's unit theorem.  The diagonal $3$-adic
logarithm map to $\prod_{v\mid3}F_v$ is $G$-equivariant.
A nonzero kernel after tensoring with $\mathbb Q_3$ would,
after scalar extension, contain a full character line.
That line is defined using algebraic character values,
so it would give a nonzero algebraic-coefficient relation
among the logarithms of multiplicatively independent
global units, at any fixed embedding into
$\overline{\mathbb Q}_3$.

Those logarithms are linearly independent over $\mathbb Q$:
an integral relation would make a product of independent
units a local root of unity, hence a global root of unity.
The Baker--Brumer theorem makes them independent over the
algebraic closure of $\mathbb Q$ as well, a contradiction.
This proves injectivity on the free part after tensoring
with $\mathbb Q_3$.  Any integral kernel is therefore
torsion.  The remaining finite $3$-power roots of unity
inject into every local multiplicative $3$-completion,
so the integral kernel is also zero.

The transcendence input is the Baker--Brumer theorem in
\citet[Theorem, p.~11-01; proof, pp.~11-02--11-08]{Rousseau1969}.
\end{proof}

\begin{lemma}[A single prime in the real subfield]
\label{lem:root-real-q-units}
Let $F/\mathbb Q$ be a finite abelian extension, let $F^+$ be its maximal
real subfield, and let $q\ne3$ be a rational prime having exactly one prime
in $F^+$.  Let $S_F$ consist of the primes above $3$ and $q$.  Then
\[
 H^1(G_{F,S_F},\mathbb Z_3(1))
 \longrightarrow\prod_{v\mid3}H^1(F_v,\mathbb Z_3(1))
\]
is injective.
\end{lemma}
\begin{proof}
Integral Kummer theory identifies the source with the $3$-completion
of the finitely generated $S_F$-unit group.  The inverse-limit ideal-class
term vanishes because the ideal class group is finite.  A class in the
kernel has zero valuation at every prime above $3$.  Exactness of
completion for finitely generated abelian groups places it in the
completed group $E_q(F)=\mathcal O_F[1/q]^\times\widehat{\ }_3$.

If $F$ is totally real, the unique prime $\mathfrak q\mid q$ has a
valuation $a\in\mathbb Z_3$.  The norm to $\mathbb Q$, followed by
localization at $3$ and the logarithm, gives
$f_{\mathfrak q}a\log_3(q)=0$.  Since $\log_3(q)\ne0$, $a=0$.
The class is then an ordinary unit class and abelian Leopoldt makes it
zero, including its integral torsion part.

Otherwise $F$ is a CM field.  Write $c$ for complex conjugation.
The integral projectors $(1+c)/2$ and $(1-c)/2$ split $E_q(F)$, since
$2$ is a unit in $\mathbb Z_3$.  Norm and inclusion identify the plus
part with $E_q(F^+)$: norm composed with inclusion is multiplication by two.
The localization kernel on this part is zero by the preceding real-field
argument, applied to $F^+$.  Here norm at $3$ uses all the places above
each place of $F^+$, as does the localization map in the statement.

It remains to consider the minus part.  Modulo roots of unity the minus
part of the ordinary unit group has rank zero.  Indeed, for every
ordinary unit $u$ the element $u/\bar u$ is a root of unity, by
Kronecker's theorem.  The valuation lattice at $q$ has minus rank at
most one: the unique prime of $F^+$ either stays a single prime in $F$
or splits into one conjugate pair.  The ideal class group is finite,
so the valuation image has finite index in the permitted lattice.
Consequently the minus part of $E_q(F)$ modulo finite torsion has
$\mathbb Z_3$-rank zero or one.

In the rank-one case choose a nonzero algebraic anti-invariant
$q$-unit after tensoring with $\mathbb Q$.  Concretely, if
$(a)=\mathfrak q^h$, then $u=a/\bar a$ has nonzero divisor
$h(\mathfrak q-\bar{\mathfrak q})$ and spans the required rational
line.  At every place above $3$, it is a local unit and its $3$-adic
logarithm is nonzero: a local unit with logarithm zero is a root of
unity, and an embedding of $F$ cannot turn this algebraic non-root
of unity into a root of unity.  Localization is therefore injective
on that line after tensoring with $\mathbb Q_3$.  Its integral kernel
can only be torsion.  The finite global $3$-power roots of unity inject
into local multiplicative $3$-completion.  This kills the remaining
kernel and proves the lemma.
\end{proof}

\begin{proposition}[Exact range of the two-prime boundary argument]
\label{prop:root-real-boundary-range}
Let $n=3^r q^s$, where $r,s\ge1$ and $q\ne3$ is an odd prime.
Choose distinct rational cusps $P,R$ and split the boundary Tate module
as $B_n=I_P\oplus M_n$ by eliminating $e_R$ from the diagonal boundary
relation. For $S=\{3,q\}$, the map
\[
 H^1(G_{\mathbb Q,S},M_n)\longrightarrow H^1(\mathbb Q_3,M_n)
\]
is injective if and only if
\[
 r=1\quad\hbox{or}\quad q\equiv2,4,5,7\pmod9.
 \tag{R}\label{eq:root-real-prime-condition}
\]
\end{proposition}
\begin{proof}
The boundary consists of three copies of $\mu_n$. Removing the two
rational basis vectors gives an integral equivariant decomposition
\[
 M_n\simeq\mathbb Z_3(1)\oplus
 \bigoplus_{\substack{d\mid n\\d>1}}
 \left(\operatorname{Ind}_{G_{F_d}}^{G_{\mathbb Q}}
       \mathbb Z_3(1)\right)^{\oplus3},
 \qquad F_d=\mathbb Q(\zeta_d).
\]
No division by a Galois-group order is involved. Shapiro's lemma reduces
the question to $S$-unit localization in these fields and in $\mathbb Q$;
the local target is the product over all places above three.
Write $d=3^a q^b$.  The decomposition subgroup at $q$ is
\[
 \langle q\bmod3^a\rangle\times(\mathbb Z/q^b\mathbb Z)^\times.
\]
Its image in $\operatorname{Gal}(F_d^+/\mathbb Q)$ is the whole group
exactly when $q$ and $-1$ generate $(\mathbb Z/3^a\mathbb Z)^\times$.
For $a\le1$ this always holds.  For every $2\le a\le r$ it holds
exactly when $v_3(q^2-1)=1$, equivalently when $q\equiv2,4,5,7\pmod9$.
Under (R), every $F_d^+$ has a unique prime above $q$, so
Lemma~\ref{lem:root-real-q-units} proves the desired injectivity.

Conversely suppose $r\ge2$ and $q\equiv1$ or $8\pmod9$.
Set $F=\mathbb Q(\zeta_{3^r})$ and $E=F^+$, of degree $g$.
There are $t\ge3$ primes of $E$ above $q$: its decomposition group
has index divisible by three in the cyclic real cyclotomic group.
The $q$-unit group of $E$ has rank $g-1+t$, whereas the product of
local unit groups at all its places above $3$ has $\mathbb Z_3$-rank $g$.
All $q$-units have zero $3$-valuations.  Thus the localization kernel
on their $3$-completion has rank at least $t-1>0$.
Restriction from $E$ to $F$ is injective on integral first cohomology,
because corestriction composed with restriction is multiplication by two.
The resulting nonzero kernel in $F$ occurs in the induced summands of
$M_n$ corresponding to $d=3^r$.  Hence the displayed map is not
injective.  This is a cohomological obstruction; it does not produce
a Selmer section with the additional local geometric conditions.
\end{proof}

\begin{theorem}[A compact Selmer criterion for mixed abelian cuspidality]
\label{thm:root-real-mixed-criterion}
For $n=3^r q^s$ as above satisfying
\eqref{eq:root-real-prime-condition}, suppose
\[
 \ker\!\left(
 \operatorname{Sel}(T_3J_n)
 \longrightarrow H^1(\mathbb Q_3,T_3J_n)
 \right)=0.
 \tag{$\mathrm C_n$}
\]
Then every Selmer section in the geometrically
pro-$3$ abelian quotient of $U_n/\mathbb Q$ is
globally cuspidal.  Its center and parameter are
unique.  Thus the Selmer set consists of the
three rational cusp packets, each parametrized
by $H^1(\mathbb Q,\mathbb Z_3(1))$.

A sufficient arithmetic condition for
$(\mathrm C_n)$ is rank $J_n(\mathbb Q)=0$
and $T_3\Sha(J_n/\mathbb Q)=0$.
The conjunction of these sufficient conditions is not asserted here
for any mixed-exponent Jacobian.
\end{theorem}
\begin{proof}
The compact cubic detector identifies the local
center at $3$ with one of the three rational
coordinate cusps $P$ of $X_n$: this uses only
the morphism $X_n\to X_3$ and the established
cubic Selmer input, not injectivity of a map
between geometric pro-$3$ groups.

Use a standard tangent at $P$ as origin.
The compact class $a$ of the section belongs
to $\operatorname{Sel}(T_3J_n)$ and has
$a_3=0$.  Hypothesis $(\mathrm C_n)$ gives
$a=0$ globally.

At every $p\notin S$, the pair $(X_n,D_n)$
has good reduction.  If a local representative
is a cusp $Q$, or an ordinary point reducing
to a cusp $Q$, its compact Kummer class is
that of $Q$.  In the ordinary case this uses
the pro-$p$, uniquely $3$-divisible formal
reduction kernel and the unique lift of the
finite \'etale boundary point.  Since this
class equals that of $P$, the $3$-primary
Kummer class of $Q-P$ is zero.
If $Q\ne P$, Lemma~\ref{lem:lm-exact-cusp-order}
makes its order exactly $n$, divisible by $3$;
its nonzero $3$-primary torsion component
cannot vanish in local $3$-completion.
Thus $Q=P$.

The zero compact class has a unique lift
$b\in H^1(\mathbb Q,B_n)$, because
$(T_3J_n)^{G_{\mathbb Q}}=0$.  Write
$b=(\beta,b_M)$ using the splitting of
Proposition~\ref{prop:root-real-boundary-range}.
The exact local disc inertia calculation
shows that $b_M$ is unramified at every
$p\notin S$: nonzero geometric inertia
can occur only along $I_P$.
At $3$, local cuspidality at $P$ and
$(T_3J_n)^{G_{\mathbb Q_3}}=0$ give
$(b_M)_3=0$.
No unramifiedness at the bad prime $q$
is claimed or needed.
Proposition~\ref{prop:root-real-boundary-range} gives
$b_M=0$ globally.

Hence the section is the cusp section
$t_{P,\beta}$.  The direct summand
$I_P\subset B_n$ and the vanishing of
compact Tate-module fixed points make
its parameter unique.  The cubic
local detector makes its center unique.
Every global cusp section is Selmer.

Finally, under the stated sufficient
arithmetic condition,
$\operatorname{Sel}(T_3J_n)$ is the finite
group $J_n(\mathbb Q)[3^\infty]$.
This group injects into local
$3$-completion at $3$, so
$(\mathrm C_n)$ follows.
\end{proof}

\begin{proposition}[The additional gates for nonabelian quotients]
\label{prop:lm-real-mixed-higher-criterion}
Retain \eqref{eq:root-real-prime-condition} and $(\mathrm C_n)$,
and let $R$ be a
$G_{\mathbb Q}$-stable quotient of $N_n$
retaining $N_n^{\mathrm{ab}}$, with finite
free ordinary lower-central factors
$L_m(R)=\gamma_mR/\gamma_{m+1}R$.
For every $m\ge2$ assume
\[
 L_m(R)^{G_{\mathbb Q_3}}=0,\qquad
 \ker\!\left(
 H^1(G_{\mathbb Q,S},L_m(R))
 \longrightarrow H^1(\mathbb Q_3,L_m(R))
 \right)=0.
 \tag{$\mathrm G_m$}
\]
Then every Selmer section in this quotient
is globally cuspidal.  For a finite
lower-central quotient only its nonzero
grades need be checked.  For a finite-derived
quotient all its nonzero grades are involved.
These additional cohomological conditions
are not proved by the $S$-unit argument.
\end{proposition}
\begin{proof}
The preceding theorem supplies the global
abelian cusp section $t=t_{P,\beta}$.
At $3$ the full local representative is
cuspidal at $P$, with parameter $\beta_3$,
by the same abelian parameter injection.
At $p\notin S$ the compact-center argument
and local inertia formula apply.  If
$v_p(\beta)\ne0$, the full free pro-$3$
normalizer argument makes the local
representative a cusp section conjugate
to $t_p$.  If $v_p(\beta)=0$, both
whole geometric inertias are zero.

For good $p\notin S$, positive Frobenius
weights give $L_m(R)^{G_{\mathbb Q_p}}=0$.
At $3$, this is automatic for $L_1$,
which is an extension of the boundary
Tate module by $T_3J_n$, and is explicitly
assumed for the higher grades.
Thus the quotient $R/\gamma_mR$ has no
fixed points under the action defined
by $t$ at the places needed here.
Inner twisting by $t$ leaves the graded
actions unchanged.

The central induction in the proof of Theorem~\ref{thm:lm-power-all-derived} therefore
makes the next difference cocycle
unramified outside $S$ and locally
a coboundary at $3$.  The original
tangent action on the grade is also
unramified outside $S$, so this
cocycle belongs to
$H^1(G_{\mathbb Q,S},L_m(R))$.
The second part of $(\mathrm G_m)$
makes it a global coboundary.
No condition is imposed on its
inertia at $q$, which is retained
inside the allowed set $S$.
The induction and compactness of
nested conjugator sets identify
the two sections in $R$.
\end{proof}

\begin{remark}[The reduction and its limits]\label{rem:real-mixed-criterion-limits}
This proves that, for the stated two-prime
mixed exponents, compact Selmer localization
injectivity is sufficient for the complete
\emph{abelian} classification.
The boundary obstruction has been removed
by a proved integral localization theorem,
after the other local geometric conditions
restrict its support to $S$.
This does not claim that the unrestricted
relative one-place cohomology kernel is zero;
the previously exhibited nonzero classes
in that kernel need not have the support
and geometry established in this argument.

The criterion is sufficient, not asserted
necessary.  In particular, failure of
$(\mathrm C_n)$ would not construct a
noncuspidal Selmer section.
Nor does $(\mathrm C_n)$ alone prove the
full or metabelian classification: the
higher-grade conditions above remain.
This proves Theorem~\ref{thm:realconditionalmixedintro}.
The proof does not transfer across a
prime-to-$3$ cover by assuming an
arithmetic splitting or a derived-quotient
embedding.
\end{remark}

